\section{Marco metodológico y conceptual}
\label{ch:marco_conceptual}

\section{Vulnerabilidad socioeconómica}
\label{sec:vulnerabilidad}

La vulnerabilidad socioeconómica se entiende aquí como la \textbf{propensión de una población territorial a experimentar privaciones en múltiples dimensiones del bienestar} (salud, educación, trabajo, vivienda, servicios públicos, ingreso), de modo que la ocurrencia de eventos adversos o la ausencia de capacidades de respuesta generan deterioros sostenidos en las condiciones de vida \citep{alkire2011multidimensional, cepal2016horizontes}.

A diferencia de la pobreza monetaria unidimensional, el enfoque de vulnerabilidad adopta una perspectiva \textbf{multidimensional, relacional y territorial}: las privaciones se manifiestan en combinaciones específicas según el contexto local, y la capacidad de respuesta depende de activos individuales, colectivos e institucionales disponibles en el territorio \citep{sen1999development, rodriguez2010vulnerabilidad}.

\section{Medición multidimensional}
\label{sec:medicion_multidimensional}

La medición de la vulnerabilidad socioeconómica requiere:

\begin{enumerate}
    \item \textbf{Selección de dimensiones y indicadores} teóricamente fundados y empíricamente observables.
    \item \textbf{Construcción de indicadores comparables} a nivel municipal (tasas, proporciones, índices normalizados).
    \item \textbf{Agregación} que respete la multidimensionalidad (evitando compensabilidad indebida entre dimensiones).
    \item \textbf{Validación} interna (consistencia estadística) y externa (con criterios independientes).
\end{enumerate}

En el contexto colombiano, el \textbf{Índice de Pobreza Multidimensional (IPM)} del DANE constituye el referente oficial \citep{dane2019ipm}. El IPM 2018 incluye 15 privaciones agrupadas en 5 dimensiones (educación, niñez y juventud, trabajo, salud, vivienda y servicios públicos). Este proyecto \textbf{no replica el IPM}; construye un índice sintético propio (IVS) a partir de indicadores censales complementarios y no redundantes con el IPM, reservando al IPM la función de \textit{validación externa posterior} (Sección \ref{sec:ipm_validacion}).

\section{Indicadores socioeconómicos}
\label{sec:indicadores_conceptual}

Un indicador socioeconómico municipal debe cumplir:

\begin{itemize}
    \item \textbf{Relevancia:} Captar una faceta de vulnerabilidad (privación, carencia, rezago, exclusión).
    \item \textbf{Comparabilidad:} Mismo denominador conceptual y escala homogénea (proporción $[0,1]$ o tasa por $10^3$/ $10^4$ hab.).
    \item \textbf{Direccionalidad:} Valor alto $\Rightarrow$ mayor vulnerabilidad (o se documenta la excepción).
    \item \textbf{Independencia relativa:} Baja colinealidad con otros indicadores de la matriz; no duplicar privaciones ya capturadas por el IPM si se busca validación externa no circular.
    \item \textbf{Disponibilidad:} Observado para la totalidad (o casi totalidad) del universo municipal.
\end{itemize}

\section{Reducción de dimensionalidad y PCA}
\label{sec:pca_conceptual}

El Análisis de Componentes Principales (PCA) es una técnica de \textbf{reducción de dimensionalidad} que transforma un conjunto de variables correlacionadas en un conjunto menor de componentes no correlacionados (componentes principales), combinaciones lineales de las variables originales que explican la máxima varianza posible \citep{jolliffe2002principal, abdi2010principal}.

En el contexto de índices compuestos de vulnerabilidad, el PCA se utiliza para:

\begin{itemize}
    \item Identificar la estructura común latente entre indicadores.
    \item Sintetizar la información en pocos componentes (idealmente el primer componente, PC1).
    \item Construir un índice sintético (scores del PC1) ponderado por la estructura de covarianzas/correlaciones de los datos, no por pesos arbitrarios.
\end{itemize}

\textbf{Supuestos y precauciones críticas:}

\begin{enumerate}
    \item \textbf{Adecuación multivariada:} KMO global $>0,6$ (idealmente $>0,7$), MSA individual $>0,5$, Bartlett significativa \citep{kaiser1974little}.
    \item \textbf{Linealidad y correlaciones:} El PCA captura estructura lineal; correlaciones muy altas ($|r|>0,9$) indican redundancia; correlaciones muy bajas sugieren ausencia de factor común.
    \item \textbf{Distribución:} Asimetría y curtosis extremas pueden distorsionar componentes; se evalúan transformaciones (log, raíz cuadrada, winsorización).
    \item \textbf{Escalado:} Se estandarizan variables (media 0, desviación 1) cuando tienen unidades o rangos distintos; aquí todas son proporciones $[0,1]$.
    \item \textbf{Interpretabilidad:} El PC1 debe representar un constructo teóricamente coherente (``vulnerabilidad general''), no un artefacto matemático.
    \item \textbf{No circularidad con validación externa:} Si el IPM se usa para validar el IVS, los indicadores que replican directamente privaciones del IPM deben excluirse de la matriz PCA o tratarse con extrema cautela \citep{booysen2002overview}.
\end{enumerate}

\textbf{Estado actual:} Los diagnósticos preliminares (KMO $\approx 0,642$, Bartlett $p<0,001$, determinante $\approx 0,0154$, MSA estrato 1--2 $\approx 0,414$) son \textbf{insuficientes} para aprobar el PCA. La auditoría estadística integral está pendiente (Sección \ref{sec:auditoria_pendiente_limitaciones}).

\section{Relación futura con contratación pública (SECOP II)}
\label{sec:secop_futuro}

La integración con SECOP II es una etapa \textbf{posterior y condicionada} a la consolidación del IVS. El diseño metodológico previsto contempla:

\begin{itemize}
    \item \textbf{Unidad de análisis municipal-año:} Panel de municipios (1.122) $\times$ años (ventana exploratoria 2015--2025).
    \item \textbf{Indicadores SECOP II:} Valor total contratado, número de contratos, concentración (Herfindahl), tipos de contrato, sectores, adjudicatarios.
    \item \textbf{Denominadores:} PIB municipal, población, presupuesto municipal, o valor total de contratación nacional como referencia.
    \item \textbf{Análisis de asociación:} Correlaciones, modelos de panel, análisis de correspondencia, pruebas de hipótesis sobre focalización.
    \item \textbf{Análisis espacial:} Autocorrelación espacial (Moran's I), modelos espaciales, clústeres de vulnerabilidad-contratación, solo si el diagnóstico lo justifica.
\end{itemize}

El \texttt{crosswalk\_enlace.csv} documenta las equivalencias territoriales pendientes para este enlace futuro (códigos 94663 $\rightarrow$ 94343, hipótesis 27493 $\rightarrow$ 27615).